% !TeX root = ../papers/nc-p3.tex
\section{計量を固定した滑らかな判別三次曲線の四割線判定}
\label{sec:marked-fixed-metric-elliptic}
\begin{theorem}\label{thm:marked-fixed-metric-elliptic}
滑らかな三次曲線 $E\subset\mathbf P(V)$ と非退化な計量を位置関係も含めて指定する。
$H=\mathcal O_E(1)$, $L=H^2$ とし、$\nu_2(E)\subset\mathbf P(H^0(L)^*)=\mathbf P^5$ を考える。
計量の逆テンソルを $o\in\mathbf P^5$ とする。$M^2=L$ なる四つの次数三線束ごとに、
$H^0(L)=\mathrm{Sym}^2 H^0(M)$ による階数一点のVeronese曲面 $V_M$ を取る。
以下は同値である。
\begin{enumerate}
\item $o$ が $M\ne H$ の三曲面のいずれにも入らない。
\item $E$ が三交換子独立な有限混合印の二次鉛筆の基点を含む。
\item そのような印を、四基点が被約な射影枠となるように選べる。
\end{enumerate}
これらのとき、この計量と $E$ そのものを固定した階数10の族があり、
添字保持同型類の像は八次元、幾何学的一般点は有限な標準Koszul左右AS正則窓で $g=5$ である。
例外では、$E$ 上のどの跡零二次鉛筆の長さ四の基点にも可換二平面を生む長さ三の部分がある。
この障害は混合印からの構成に関するものである。
\end{theorem}
\begin{proof}[六次埋込みの割線]
次数三線束は平面三次埋込みを与え、そこに二次方程式はないから
$\mathrm{Sym}^2H^0(M)\to H^0(L)$ は同型である。$o$ は元の平面に対し階数3なので、
非被約な場合も含め長さ2以下の二次Veronese像の包には入らない。
長さ四の有効因子 $D$ が $o$ をその包に含む条件は、$D$ 上で消える二次式が全て
$o$ と跡直交する条件である。この四点が射影枠ならその二次式はちょうど二次鉛筆で、
平面での共通零スキームが $D$ に等しい。

まず $o$ を包に含む次数三因子を調べる。相異なる二つ $T,T'$ は交わらない。
交われば和集合の長さは5以下で、その包の次元は長さ引く1であり、
二つの包の共通部分は長さ2以下の共通部分の包になる。これは種数1のRiemann--Rochと
$o$ の階数に反する。交わらないなら、六次因子の包が予想より小さい条件は
$\mathcal O_E(T+T')=L$ である。従って三つ以上あれば全て同じ平方根 $M$ の切断である。
$T=(s)_0$ が条件を満たすことは、$o$ が $s H^0(M)$ を消すこと、
すなわち $s$ が $H^0(M)$ 上の対称形式 $o$ の核に入ることに等しい。
よって無限個の場合に限って $o\in V_M$ であり、その全てがこの核の射影直線をなす。
$o\notin V_H$ は階数3より明らかである。

次に $A\in\mathrm{Pic}^4(E)$ とし、$N=L A^{-1}$ と置く。
$D\in|A|$ が $o$ を張る条件は、その切断が積と $o$ で定まる
$H^0(A)\to H^0(N)^*$ の核に入ることである。この写像の階数は常に2である。
実際双対に核 $u$ があれば $o$ は $uH^0(A)$ を消し、長さ2の零因子 $(u)_0$ の包に
入ってしまう。従って全ての $D$ はPicard曲線上の射影直線束を動く。
上記の三因子が有限なら、一般の $A$ でこの鉛筆は基点自由である。
$x$ が基点なら $H^0(A(-x))$ の三次元空間の像が一次元以下となり、
ある非零 $u\in H^0(N)$ について $o$ は $x+(u)_0$ の包に入る。
有限個の三因子とその支持 $x$ から $N=\mathcal O_E((u)_0)$ が決まるため、
該当する $A$ は有限個である。
標数零で一般切断の零点は相異なる。三つが共線になる場合、第四点 $y$ は $A=H(y)$ で一意に定まり、
基点自由な鉛筆の一般元はそれを避ける。有限個の三因子を含むものも避けられる。
従って一般の $D=\{v_1,\ldots,v_4\}$ は射影枠で、$o$ を四つの平方の線型結合にした際の
四係数は全て非零である。

一方 $o\in V_M$ の場合、任意の $A$ に対し $A=M(x)$ となる $x$ は一意である。
上で求めた三因子の鉛筆に固定点 $x$ を加えると $|A|$ 内に条件を満たす射影直線が得られる。
すでに核の次元は2なので、これで全てである。特にどの $D$ も、$o$ を張る長さ三因子 $T$ を含む。
その局所環全体 $R=H^0(\mathcal O_T)$ は三次元である。$T$ 上単元の一次式で自明化すると、
一次式の $R$ への制限は同型である。さもなくば $T$ は一本の直線内にあり、
その平方の包には階数3の $o$ が入らない。$o$ は $T$ の包にあるからある $\tau:R\to k$ で
$\tau(ab)$ の積対合を一次式空間に引き戻したものとして表される。この対合は非退化なので、
$T$ の平方の線型包を計量で自己準同型に戻すと、$R$ の全乗法作用の空間となる。
実際その全双対は $r\mapsto\tau(sr)$ の形に尽きる。これは可換で単位を含む三次元空間である。
$D$ の二次鉛筆の直交補の跡零部分は、従って可換二平面を含み、
三交換子の独立性とは両立しない。これは重なった点の接方向にも適用される。
\end{proof}

\begin{proof}[射影枠から有限混合印へ]
計量を標準形にし、一般の $D$ の四つの代表を列とする $3\times4$ 行列 $P$ を
$PP^t=I$ となるよう取る。これは上の四係数の非零平方根を使えばよい。
$\ker P$ は非退化な直線であり、その生成元を $h^th=1$ に取る。どの三列も独立なので全 $h_i\ne0$。
$K=k^4$ に $\tau(x)=\sum h_i^2 x_i$ を入れる。$U=\ker\tau$ から
$\operatorname{diag}(h_i)$ により $h^\perp$ を経て $P$ で三空間に移ると、
符号の平方根を除いて $-\tau(xy)$ は元の計量になる。
$K$ の原始冪等元による乗法を $U$ に射影した行列は、非零倍率を除いて
$P$ の同名の列の平方である。実際 $h^\perp$ への直交射影は $P^tP$ で、
座標射影を両側から圧縮すると $P e_i e_i^t P^t$ となる。
跡零部分の包は従って $D$ の二次鉛筆に対応する $G'$ である。
定理\ref{thm:mixed-point-frobenius-algebras}の構成で $m$ を変えてもこの包は同じで、
順序付き基底だけが変わる。相手が可逆な外積三列なので $G'$ の三交換子は独立である。

一般の $m$ で辺が有限であることも確認する。上記定理の二周期式を使い、
$L^tL$ に当たる自己随伴同型を $M_0$、$n=L^ty$ と書く。
星点以外なら $x,y$ は共に非零である。$r=0$ のとき $y\parallel Lx$ なので
$B_{M_0x}x\parallel x$ が必要である。ここで $B_zx$ は $K$ の積の $U$ 成分である。
$M_0=I$ では $x^2\in\langle1,x\rangle$、すなわち四座標は相異なる二値を取る。
二つの値の組分けは有限、跡零条件により各組分けで射影点は一意なので、この軌跡は有限である。
射影的な繊維次元の上半連続性により一般の $M_0$ にも有限性が移る。
$r\ne0$ なら外積の式から $y\cdot y=0$ と $x\cdot n=0$ を得る。
従ってもう一方の式は $K$ で $(s'1+n)x=0$ となる。
$\tau(x)=0$ かつ重みが全て非零なので、$n$ の少なくとも二座標が等しい。
また $n$ は $n\cdot M_0^{-1}n=0$ の非退化二次曲線上にある。
一般の $M_0$ ではこれは六本の等値直線の交点を避け、各直線と有限個の点で交わる。
従って重複する一対が決まって $x$ はその二成分に支えられる跡零の一方向である。
残る星座標も式から決まり、符号を逆にした辺も同じ議論で有限となる。

こうして射影直線束の稠密開部分の全ての $D$ で有限混合印を得る。前節の
指定三次式への滑らかさとエタール復元から、各 $D$ で一般点が全次数正則な
十六次元の表示が得られる（$G'$ の九基底座標を含む）。
全体ではこの二次元の基点族に沿う同じ構成を用いる。二つの交代解の符号は残余群が交換し、
他の基底の開条件と混合印を通る局所成分は前節までと同じである。
一般窓の $G'$ は指定印から決まり、$D$ が復元される。計量と $E$ を共に保つ
射影群は有限である（滑らかな平面三次曲線の射影安定化群は有限）。
よって連結な残余群の次元は10、印を忘れる選択肢は有限であり、像次元は $2+16-10=8$ となる。
\end{proof}

\paragraph{例外の具体例。}
$C=V(17xyz-4(x^3+y^3+z^3))$ は滑らかである。
\[
 D(x,y,z)=\begin{pmatrix}x&2z&2y\\2z&y&2x\\2y&2x&z\end{pmatrix},
 \qquad o=[\operatorname{adj}D(1,0,0)]
\]
と置く。余因子の六つの二次式は独立なので $[\operatorname{adj}D]$ の像はVeronese曲面である。
$C$ 上で $D$ は階数2で、核点 $v$ は $D(v)(x,y,z)^t=0$ より再び $C$ 上にある。
この対応は自分自身が逆なので、余因子の曲面は $\nu_2(C)$ を含む。
元のVeroneseとは異なり $o$ は階数3である。従って曲面の線型因子から来る次数三線束は
$H$ と異なる平方根であり、この固定計量では上の例外が実際に生じる。
